% Conteúdo do Apêndice A: Bibliotecas e versões utilizadas

\begin{table}[H]
\centering
\renewcommand{\arraystretch}{1.1}
\caption{Principais bibliotecas Python e versões utilizadas.}
\label{tab:versoes}
\small
\begin{tabular}{|l|l|p{7.2cm}|}
\hline
\textbf{Biblioteca} & \textbf{Versão} & \textbf{Finalidade} \\ \hline
SQLAlchemy & 2.0.53 & Conexão e consultas ao PostgreSQL \\ \hline
psycopg2-binary & 2.9.13 & Driver do PostgreSQL e inserção em lotes (\texttt{execute\_values}) \\ \hline
PyMongo & 4.18.1 & Conexão e consultas ao MongoDB \\ \hline
Faker & 40.39.0 & Geração de dados sintéticos \\ \hline
pandas & 2.3.3 & Manipulação dos resultados \\ \hline
matplotlib & 3.10.9 & Geração dos gráficos \\ \hline
numpy & 2.2.6 & Cálculos numéricos nos gráficos \\ \hline
python-dotenv & 1.2.3 & Leitura das variáveis do \texttt{.env} \\ \hline
\texttt{time} / \texttt{statistics} & 3.10.12$^{*}$ & Medição (\texttt{perf\_counter}) e agregação dos tempos \\ \hline
\texttt{subprocess} / \texttt{threading} & 3.10.12$^{*}$ & Coleta de CPU e memória dos contêineres (cgroup, com \texttt{docker stats} como alternativa) durante a inserção dos dados \\ \hline
\end{tabular}
\vspace{0.4em}
\legend{$^{*}$ Biblioteca padrão do Python, com a versão do interpretador. Fonte: Elaborado pela autora.}
\end{table}
